\begin{center}
    \section*{\fontsize{20}{20}\selectfont Chapter 2}
\end{center}
\vspace{10mm}
\section{Background Study}

This chapter reviews earlier work on automatic plant disease detection. It starts with classic Machine Learning (ML) methods that need hand-made features. It then moves to Deep Learning (DL) and Convolutional Neural Networks (CNNs), which changed the field. We look in particular at transfer learning, small models for mobile phones and Explainable AI (XAI), which helps people trust a model. At the end, a problem analysis lists the gaps that this thesis tries to fill.

\subsection{Literature Review}

\subsubsection{Plant Pathology: Agricultural Context and Diagnostic Challenges}
Plant diseases are a big risk to food security and to the economy. The Food and Agriculture Organization (FAO) reports that plant pests and diseases reduce world crop yields by about 20--40\% every year \cite{strange2005plant}. Diseases spread fast, and many of them look alike on the leaf. This makes early detection hard, but it is very important for farming.

The usual ways of diagnosis are well described in farming studies. Most often an expert looks at the plant and checks for signs such as lesions, spots or color changes \cite{barbedo2018factors}. But studies point out problems with this: people make mistakes, there are few experts in rural areas and checking by hand takes a long time. Lab tests can identify the pathogen correctly, but they cost a lot and need special equipment \cite{fao2017future}.

Recent farming studies stress the value of precision farming. They show that acting in time, with pesticide only where it is needed, can cut crop loss and also reduce the harm from using too much chemical \cite{savary2019global}. However, large monitoring programs are hard to run, because resources are limited and remote fields lack technology. This is why there is strong interest in cheap, accurate and easy screening tools based on Artificial Intelligence \cite{lecun2015deep}.

\subsubsection{Traditional Machine Learning Approaches}
Before deep learning, researchers mostly used image processing together with classic classifiers. Color histograms, Scale Invariant Feature Transform (SIFT) and Gray-Level Co-occurrence Matrix (GLCM) were common ways to get texture and shape features from leaves \cite{pujari2016image}. For example, Arivazhagan et al. used texture features to find unhealthy areas on plants \cite{arivazhagan2013detection}. Al-Hiary et al. built a system that used K-Means clustering for segmentation and a neural network for classification \cite{alhiary2011fast}. These methods were an important start. But they depended on hand-made features, and they often failed on busy backgrounds or large datasets \cite{ngugi2021recent}.

\subsubsection{Deep Learning in Agriculture: Evolution and Applications}
Deep learning for farm images has changed a lot in the last ten years \cite{mohanty2016using}. Key studies in 2015--2016 showed that deep learning can work for plant diseases. The PlantVillage dataset gave researchers a common benchmark. With it, they showed that CNNs can learn features from leaf images layer by layer, without any hand-made features \cite{hughes2015open}.

Mohanty et al. (2016) tested AlexNet and GoogLeNet on PlantVillage. They got more than 99\% accuracy, but they also saw that accuracy dropped a lot on images with different backgrounds \cite{mohanty2016using}. Sladojevic et al. (2016) built a model that recognizes 13 types of plant disease, which showed that Deep Neural Networks (DNNs) are useful in farming \cite{sladojevic2016deep}. Later work used deeper networks. Ferentinos (2018) tried VGG16, AlexNet and GoogLeNet. VGG16 gave the best accuracy (99.53\%), but the author noted that it is very costly to run, because it has about 138 million parameters \cite{ferentinos2018deep}.

Most of these studies used lab photos with a plain background. To fix this, Shabrina et al. (2024) released a potato leaf disease dataset taken in the field \cite{shabrina2024novel}. The photos show soil, hands, plastic mulch, shadows and other leaves, so the task is much harder than PlantVillage. In this thesis we use our curated 5-class version of this outdoor field dataset (2,321 images).

\subsubsection{Transfer Learning and CNN Architectures}
In transfer learning, a model trained on a large dataset such as ImageNet is fine-tuned for a new task. It is now the normal way to deal with small datasets \cite{pan2010survey}. Three model families matter for this work.

\textit{Multi-Backbone Ensembles (ResNet50, EfficientNet and MobileNet).} He et al. (2016) introduced ResNet with ``skip connections'' to solve vanishing gradients \cite{he2016deep}, and Tan and Le (2019) introduced EfficientNet \cite{tan2019efficientnet}. Combining EfficientNetB0, MobileNetV2 and ResNet50 in one ensemble (MultiNet-V3) gives a strong heavy benchmark model.

\textit{MobileNetV2.} Sandler et al. (2018) proposed MobileNetV2 for phones and other small devices. It uses depthwise separable convolutions and inverted residuals, which cut the number of parameters and the delay by a large amount \cite{sandler2018mobilenetv2}. With careful fine-tuning, it can reach higher accuracy than much larger models on field photos.

\textit{VGG and InceptionV3.} VGG19 \cite{simonyan2015very} and InceptionV3 \cite{szegedy2016rethinking} are also widely used. InceptionV3 uses parallel multi-scale filters at $299 \times 299$ resolution, making it a strong single-backbone benchmark.

\subsubsection{Data Augmentation and Optimization Techniques}
Several methods help a model work on new images and stop it from overfitting. Data augmentation (rotation, flipping, zoom) creates more training images from the existing ones, so the model handles unseen images better \cite{shorten2019survey}. Dropout, proposed by Srivastava et al., turns off random neurons during training, so the neurons do not depend too much on each other \cite{srivastava2014dropout}. Batch Normalization \cite{ioffe2015batch} and newer activation functions such as Swish \cite{ramachandran2017swish} also help deep networks train faster and reach higher accuracy.

\subsubsection{Explainable AI (XAI) and Model Transparency}
A common complaint about deep learning in farming is that a CNN is a ``black box''. Farmers need to know why a model predicts a certain disease. Selvaraju et al. (2017) proposed Grad-CAM (Gradient-weighted Class Activation Mapping), which shows the regions of the image that the model used \cite{selvaraju2017grad}. Newer studies have started to use Grad-CAM and other interpretable ML methods to check that the model looks at leaf lesions and not at background noise \cite{ryo2022explainable}. Still, many proposed systems skip this check, so we cannot be sure they are reliable.

\subsubsection{Large Language Models in Agriculture}
Large Language Models (LLMs) such as GPT and Gemini can write explanations in normal language and answer questions \cite{google2026gemini}. In farming, they can turn a technical prediction into advice that a farmer understands. But LLMs can give wrong or made-up information. So their output must be limited by careful prompts and checked by experts. In this work the LLM does not make the diagnosis. It only explains the class that the CNN predicted.

Table~\ref{tab:lit_summary} sums up the main studies above. It compares their datasets, methods, results and the gaps that remain.

\begin{table}[htbp]
\caption{Summary of Key Literature in Plant Disease Detection}
\label{tab:lit_summary}
\begin{center}
\scriptsize
\begin{tabularx}{\textwidth}{| p{2.2cm} | p{2.4cm} | p{2.7cm} | p{2.3cm} | X |}
\hline
\textbf{Study} & \textbf{Dataset / Setting} & \textbf{Approach / Models} & \textbf{Reported Outcome} & \textbf{Key Gap Addressed in This Work} \\
\hline
Al-Hiary et al. (2011) \cite{alhiary2011fast} & Leaf images (controlled) & K-Means + Neural Networks & Good segmentation baseline & Uses hand-made features; sensitive to light and clutter. \\
\hline
Arivazhagan et al. (2013) \cite{arivazhagan2013detection} & Leaf texture samples & Color/GLCM texture + SVM & Accurate on small chosen sets & Manual features do not work on busy backgrounds. \\
\hline
Mohanty et al. (2016) \cite{mohanty2016using} & PlantVillage (lab plain background) & AlexNet, GoogLeNet & $>99\%$ on lab split & Large drop when tested on field or other backgrounds. \\
\hline
Sladojevic et al. (2016) \cite{sladojevic2016deep} & 13 plant disease classes & Deep CNN (CaffeNet) & Working end-to-end CNN & Needs many resources; no XAI and no advice for farmers. \\
\hline
Ferentinos (2018) \cite{ferentinos2018deep} & PlantVillage (58 classes) & VGG16, AlexNet, GoogLeNet & 99.53\% (VGG16) & $\sim 138$M parameters; too heavy for phones and edge devices. \\
\hline
Too et al. (2019) \cite{too2019comparative} & PlantVillage (fine-tuning) & VGG, ResNet, DenseNet, Inception & DenseNet-121 highest accuracy & Lab data only; no field backgrounds and no LLM advice. \\
\hline
Ryo (2022) \cite{ryo2022explainable} & Agricultural data analysis (XAI review) & Post-hoc XAI \& interpretable ML & Reviews XAI methods in farming & Shows farmers lack trust; heatmaps must be checked on lesions. \\
\hline
Shabrina et al. (2024) \cite{shabrina2024novel} & Potato Leaf (3,076 field images, 7 classes) & Benchmark field dataset & Realistic field benchmark & Very uneven classes (11:1) and busy backgrounds; no small model + XAI + LLM pipeline. \\
\hline
\end{tabularx}
\end{center}
\end{table}

\subsection{Problem Analysis}
From the review, we found three main problems that led to this thesis.

\textit{Too much computation.} Top models such as VGG16 \cite{ferentinos2018deep} and ResNet50 \cite{he2016deep} are accurate but heavy. They need a lot of memory and processing power. The cheap smartphones used by farmers in developing countries often do not have this.

\textit{Accuracy or speed.} Small models such as MobileNet solve the size problem. But they sometimes miss fine texture details of some diseases that deeper networks can see \cite{kc2019depthwise, too2019comparative}. We need a model made for this task that is as fast as MobileNet and as precise as ResNet.

\textit{No transparency.} Most systems only give a short text output, such as ``Healthy'' or ``Diseased''. Without any visual proof or explanation, people do not trust them. XAI is needed to check that high accuracy is real and does not come from bias in the data \cite{ryo2022explainable, amara2017deep}.

\noindent \textbf{Conclusion of Analysis:} To solve these problems, this research uses four parts: (1) a compact fine-tuned MobileNetV2 model (8.64~MB) to keep the computing cost low, (2) a comparison with MultiNet-V3 Ensemble and InceptionV3 to measure its performance, (3) Grad-CAM to make the decisions visible, and (4) an LLM advice module to make the output useful for farmers.
